
\subsubsection{2.1 Weight-Space Self-Supervised Learning}

Learning representations from neural network weights --- treating checkpoints as data --- has attracted growing attention. Hyper-Representations [Schurholt2021Hyper] first demonstrated that an autoencoder trained on populations of neural network weights produces latent codes that recover hyperparameters, accuracy, and generalization gap without task labels. SANE [Schurholt2024Towards] extended this to a hierarchical transformer architecture with contrastive SSL, achieving $R^{2}=0.72$ on heterogeneous model zoo property prediction within its intended same-architecture setting. Both methods use flat tokenization: weight matrices are chunked into fixed-size tokens, processed positionally, and aggregated --- an approach that implicitly assumes weight token positions are comparable across networks.

This assumption holds when training and test architectures share the same layer structure. When architectures differ substantially (e.g., MLP/CNN vs. ViT), flat tokenization loses its grounding: the positional identity of tokens has no consistent meaning across the architectural boundary. Our work measures the consequence empirically: SANE's representations collapse to near-constant codes for cross-architecture inputs ($std\approx 0.0014)$, and property prediction falls to $R^{2}=0.072$. This is a cross-architecture distribution shift failure, not a failure of SANE in its intended domain.

\subsubsection{2.2 Equivariant Graph Neural Networks for Weight Encoding}

A parallel research strand represents neural networks as computation graphs and uses equivariant GNNs as encoders. Navon et al. [Navon2023Equivariant] characterized all permutation-equivariant and permutation-invariant linear layers over weight spaces, providing the theoretical foundation. Neural Functional Transformers [Zhou2023Permutation] and NF-Layers / Neural Graphs [Kofinas2024Graph] demonstrated that graph-based encoders achieve strong supervised property prediction within a fixed architecture family. Critically, Kofinas et al. [Kofinas2024Graph] \citep{ICLR2024} showed that graph-based encoders can generalize across architecture families in the supervised setting, achieving $R^{2}=0.71$ on held-out CNNs after training on MLPs. ScaleGMN [Kalogeropoulos2024Scale] (NeurIPS 2024 Oral) extended this to include scale equivariance, showing +8-12 $R^2$ points over permutation-only baselines in supervised same-architecture settings.

Our work operates in two ways that differ from this prior work: (1) we use the self-supervised setting, without property labels during encoder training; and (2) we target cross-architecture transfer to ViTs with LayerNorm, which changes the benefit profile of scale equivariance relative to what [Kalogeropoulos2024Scale] demonstrated in the supervised same-architecture setting.

\subsubsection{2.3 Cross-Architecture Weight Representation Transfer}

Ballerini et al. [Ballerini2025Weight] applied contrastive SSL to diverse neural radiance field architectures and demonstrated cross-architecture property transfer without target-domain training data. Our work applies a similar paradigm to general-purpose neural architectures (MLPs, CNNs, ViTs) --- a broader and more architecturally diverse setting. The ViT Model Zoo [Falk2025ModelZoo] provides real ViT-S/16 ImageNet checkpoints with accuracy labels and serves as our evaluation dataset.

\subsubsection{2.4 Normalization Layers and Symmetry Groups}

The interaction between weight-space symmetry groups and neural network normalization layers has received theoretical attention. Scale equivariance addresses the functional equivalence $\theta$ ~ $c\theta$ for scalar c > 0 (valid when activations are scale-homogeneous, as in ReLU networks). We hypothesize this equivalence breaks down in the presence of LayerNorm, which normalizes each neuron's activations to zero mean and unit variance, effectively fixing the activation scale at inference. If this gauge-fixing argument is correct, scale equivariance encodes an incorrect inductive bias for ViT weight encoders. Our empirical finding --- scale+permutation equivariance underperforms permutation-only --- is consistent with this hypothesis. An unverified preprint [Anonymous2025LayerNorm; arXiv:2510.08300] appears to develop a related theoretical claim; we flag this citation as unverified at the time of submission and hedge our claims accordingly.

